\begin{tikzpicture}[
    font=\small,
    flow/.style={
        draw=black,
        -Latex,
        line width=0.8pt
    },
    dashedflow/.style={
        draw=black!75,
        -Latex,
        dashed,
        line width=0.75pt
    },
    panel/.style={
        draw=black,
        rounded corners=4pt,
        line width=0.8pt
    },
    titlebox/.style={
        draw=black,
        rounded corners=4pt,
        fill=blue!5,
        minimum height=0.72cm,
        align=center
    },
    stepbox/.style={
        draw=black,
        rounded corners=4pt,
        fill=gray!5,
        minimum width=3.8cm,
        minimum height=1.1cm,
        align=center,
        line width=0.8pt
    },
    actorbox/.style={
        draw=black,
        rounded corners=4pt,
        fill=blue!8,
        minimum width=4.1cm,
        minimum height=1.05cm,
        align=center,
        line width=0.8pt
    },
    criticbox/.style={
        draw=black,
        rounded corners=4pt,
        fill=green!8,
        minimum width=4.1cm,
        minimum height=1.05cm,
        align=center,
        line width=0.8pt
    },
    envbox/.style={
        draw=black,
        rounded corners=4pt,
        fill=orange!10,
        minimum width=4.1cm,
        minimum height=1.05cm,
        align=center,
        line width=0.8pt
    },
    notebox/.style={
        draw=black,
        rounded corners=4pt,
        fill=yellow!10,
        minimum width=5.2cm,
        minimum height=1.0cm,
        align=center,
        line width=0.8pt
    },
    comparebox/.style={
        draw=black,
        rounded corners=4pt,
        fill=gray!8,
        minimum width=5.4cm,
        minimum height=1.8cm,
        align=left,
        line width=0.8pt
    }
]

% =========================
% Outer panel
% =========================
\draw[panel, fill=white] (0.2,0.3) rectangle (18.7,11.2);
\node[titlebox, minimum width=5.4cm] at (9.45,11.7) {PPO 的策略更新与样本重复利用过程};

% =========================
% Main steps
% =========================
\node[actorbox] (s1) at (3.0,9.7) {
    第 1 步\\
    固定当前策略\\
    $\theta \rightarrow \theta_{\mathrm{old}}$
};

\node[envbox] (s2) at (9.45,9.7) {
    第 2 步\\
    使用旧策略与环境交互\\
    采集状态、动作、奖励数据
};

\node[criticbox] (s3) at (15.8,9.7) {
    第 3 步\\
    Critic 计算价值估计\\
    得到回报和优势 $\widehat{A}_t$
};

\node[stepbox] (s4) at (3.0,6.9) {
    第 4 步\\
    将采集到的数据\\
    划分为多个小批次
};

\node[stepbox] (s5) at (9.45,6.9) {
    第 5 步\\
    基于裁剪目标\\
    更新 Actor
};

\node[criticbox] (s6) at (15.8,6.9) {
    第 6 步\\
    更新 Critic\\
    提高价值估计精度
};

\node[notebox] (s7) at (9.45,4.4) {
    对同一批数据重复训练 $K$ 轮\\
    每一轮都使用裁剪目标限制策略变化
};

\node[actorbox] (s8) at (9.45,1.9) {
    第 7 步\\
    完成本轮更新后\\
    用新策略重新采样数据
};

% =========================
% Main arrows
% =========================
\draw[flow] (s1.east) -- (s2.west);
\draw[flow] (s2.east) -- (s3.west);

\draw[flow] (s3.south west) -- (s4.north east);
\draw[flow] (s4.east) -- (s5.west);
\draw[flow] (s5.east) -- (s6.west);

\draw[flow] (s6.south west) -- (s7.north east);
\draw[flow] (s7.south) -- (s8.north);

% =========================
% Reuse loop
% =========================
\draw[dashedflow]
    (s7.west)
    -- node[midway, above,xshift=5mm] {\footnotesize 重复 $K$ 轮} (s4.south);

% =========================
% Comparison note
% =========================
% \node[comparebox] at (3.5,2.1) {
% \textbf{与 A2C 的区别：}\\
% A2C：采集一批数据后通常只更新一次\\
% PPO：采集一批数据后可划分小批次并重复更新多轮
% };

% =========================
% Small note
% =========================
% \node[comparebox, minimum width=5.0cm, minimum height=1.6cm] at (15.1,2.1) {
% \textbf{核心原因：}\\
% 裁剪目标限制单次策略更新幅度\\
% 因此同一批样本可以在一定范围内重复利用
% };

\end{tikzpicture}